\section{Threats to Validity}
\label{sec:threats}

The threats taxonomy follows Wohlin et al.~\cite{wohlin2012experimentation}.

\textbf{Conclusion validity.} The sweep comprises 15{,}000 (cell, seed,
window, method) observations summarised with 95\% percentile bootstrap
confidence intervals (10{,}000 resamples). All decay slopes are
strictly negative with bootstrap CIs that exclude zero. Spearman
correlations for H1 and H2 are computed on 5- and 4-point cell-mean
slope series respectively; small sample size limits the precision of
the $\rho$ estimates, and we report the point values rather than asymptotic
$p$-values. The H1/H2 rejections are not artefacts of low statistical
power: the observed $\rho$ values fall far short of the pre-registered
$\pm 0.8$ thresholds and frequently carry the wrong sign.

\textbf{Internal validity.} Three concerns.

\textit{(i) Selection-policy coupling.} All methods are evaluated on
the fleet trajectory generated by HygienePrio-full's top-$K$ selections
(\S5). At high $K$, this strongly determines which CVEs are drained
from the backlog. EPSS-only's measured decay rate therefore mixes
intrinsic EPSS compositional decay with the effect of being evaluated
against an evolved-by-HygienePrio backlog. A fully factorial
``each method drives its own trajectory'' sweep would separate these
effects, at the cost of non-identifiable cross-method comparison
between methods that produced different fleet trajectories. We discuss
this as future work and bound the present claims to the
HygienePrio-driven evolution regime.

\textit{(ii) Threshold drift.} The HRS 75th-percentile ground-truth
threshold is recomputed each window. At high $K$, this threshold
moves rapidly because the high-HRS host pool depletes. P@50 is
therefore measured against a moving positive set whose composition
co-evolves with capacity, particularly at $K \in \{100, 200\}$. A
fixed Window-1 threshold was rejected during pre-registration because
it would produce trivially shrinking positive counts that make
per-window P@50 incomparable; the chosen design preserves cross-window
comparability at the cost of entangling capacity with denominator drift.

\textit{(iii) Synthetic-fleet structural priors.} EPSS and HRS
distributions, patch-lag dynamics, KEV prevalence, and AD-exposure
structure are all drawn from the EEHDA fixture distributions inherited
from HygienePrio~\cite{paper4hygieneprio}. Real fleet distributions
may differ in ways that change the cells at which collapse occurs.
The qualitative pattern --- universal EPSS decay, capacity-dependent
HygienePrio collapse, regime-robust per-pair dominance --- is the
result of compositional dynamics that are unlikely to be artifacts of
the specific distribution choices, but the exact cell boundaries
($K = 100$ vs.\ $K = 200$ for the HygienePrio inflection) are
distributional.

\textbf{Construct validity.} P@50 with binary relevance is informative
when the positive set is large; under heavy decay, the positive count
shrinks and P@50 becomes noisier. The decay-slope metric partially
compensates by aggregating across windows. A complementary
remediation-latency metric integrating scheduling delay against
exploitation-event arrival was not in scope for this sweep
(consistent with the deferral in~\cite{paper5temporal}).

\textbf{External validity.} The 14-day window, 0.15 patch-debt
reduction, 0.02 EPSS random-walk, and (8\%, 12\%) KEV/EPSS prevalence
calibrations follow the priors of~\cite{paper4hygieneprio,paper5temporal}. Organisations with
substantially different patch-velocity heterogeneity, scanner coverage,
or AD topology may sit at $(K/\lambda)$ ratios outside the studied
grid. The qualitative finding --- that both methods exhibit
capacity-dependent absolute collapse while preserving HygienePrio's
per-pair advantage --- is more generalisable than any specific cell-level
number; no claim is made about generalisation to real fleets in the
absence of external validation on real exploitation outcome data.
